\documentclass[../../../include/open-logic-section]{subfiles}
%READERNOTE{SOL6-A224: आरंभीच्या टप्प्याला आधीचे टप्पे नसतात. संचयी मांडणीमुळे रसेलचा अमर्याद संच तयार होत नाही हा युक्तिवाद आहे; कोणत्याही औपचारिक संच उपपत्तीच्या पूर्ण सुसंगततेची सिद्धता नाही.}

\begin{document}

\olfileid{sth}{story}{approach}
\olsection{संचयी-क्रमिक पद्धत}

संचांचे टप्प्यांमध्ये
स्तरीकरण कसे
करायचे, याचे
थोडे अधिक
तपशीलवार
वर्णन असे:
\begin{quote}
  संच \emph{टप्प्यांमध्ये}
  तयार होतात.
  प्रत्येक टप्पा~$S$
  यासाठी, $S$ च्या \emph{आधी}
  कोणते टप्पे येतात, हे निश्चित असते.
  $S$ टप्प्यावर,
  $S$ पूर्वीच्या
  टप्प्यांवर
  तयार झालेले
  संच घेऊन
  बनलेल्या
  प्रत्येक
  संग्रहाचा
  संच तयार
  होतो.
  टप्प्यांवर
  तयार झालेले
  संच सोडून
  इतर संच
  नाहीत.
  \citep[p.~323]{Shoenfield:AST}
\end{quote}
ही \emph{संचांची
संचयी-क्रमिक
संकल्पना}
याची रूपरेषा
आहे. \olref[sth][][]{part}
मध्ये मांडल्या
जाणाऱ्या
औपचारिक
संच उपपत्तीला
ती आधार
देईल.

थोडा अधिक
तपशील पाहू.
शोनफिल्ड
यांच्या
वर्णनानुसार,
प्रत्येक
टप्प्यावर
मागील
टप्प्यांपासून
आपल्याला
उपलब्ध
असलेल्या
संचांपासून
नवे संच
तयार होतात.
म्हणून
आरंभीच्या
टप्प्यावर,
म्हणजे
टप्पा~$0$
येथे,
\emph{आधीचा}
टप्पाच
नसतो;
\emph{मग तर}
पूर्वीच्या
टप्प्यांतून
आपल्याला
संचही
उपलब्ध
नसतात.\footnote{
  पहिला
  टप्पा
  \emph{असतो}
  असे का
  गृहीत
  धरावे?
  \olref[z][story]{sec}
  मधील
  \stagesord{}
  जवळची
  तळटीप
  पाहा.}
त्यामुळे
आपण
एकच संच
तयार करतो:
एकही
!!{element}s
नसलेला
रिक्त संच
$\emptyset$.
टप्पा~$1$
येथे
मागील
टप्प्यांपासून
बरोबर
एक संच
उपलब्ध
असतो,
म्हणून
फक्त एक
नवा संच
$\{\emptyset\}$
तयार होतो.
टप्पा~$2$
येथे मागील
टप्प्यांपासून
दोन संच
उपलब्ध
असतात;
आपण दोन
नवे संच
$\{\{\emptyset\}\}$
आणि
$\{\emptyset,
\{\emptyset\}\}$
तयार करतो.
टप्पा~$3$
येथे मागील
टप्प्यांपासून
चार संच
उपलब्ध
असतात;
म्हणून
बारा
नवे संच
तयार
होतात\ldots.
त्यामुळे
संचांचे
संचयी-क्रमिक
चित्र
साधारण
असे
दिसेल
(अंक टप्पे
दर्शवतात):
\begin{center}
  \begin{tikzpicture}[scale=0.6]
  \tikzset{cut_here/.style={densely dotted}}
  \draw (-4,6) -- (-1, 0)-- (2,6);
  \draw[cut_here] (-5,8)--(-4,6);
  \draw[cut_here] (2,6)--(3,8);
  \draw(-1.5, 1)--(-0.5, 1);
  \draw(-2, 2)--(0, 2);
  \draw(-2.5, 3)--(0.5,3);
  \draw(-3, 4)--(1, 4);
  \draw(-3.5, 5)--(1.5, 5);
  \draw(-4, 6)--(2, 6);
  \node[label] at (0, 0) {\small 0};
  \node[label] at (0.5, 1) {\small 1};
  \node[label] at (1, 2) {\small 2};
  \node[label] at (1.5, 3) {\small 3};
  \node[label] at (2, 4) {\small 4};
  \node[label] at (2.5, 5) {\small 5};
  \node[label] at (3, 6) {\small 6};
  \end{tikzpicture}
\end{center}
मग ही मांडणी
का स्वीकारावी?

एक कारण:
ती चांगली
आणि हाताळायला
सुलभ मांडणी
आहे. सर्वांत
उघड वाटणारी
मांडणी,
म्हणजे
अनौपचारिक
गुणधर्माधारित
संचरचना,
अपयशी ठरल्यावर
आपल्याला
एखादी चांगली
मांडणी हवीच.

पण ही मांडणी \emph{फक्त} सुलभ नाही.
ती रसेलच्या विरोधापत्तीला कारणीभूत असलेली
अमर्याद संचरचना का मानू नये हे स्पष्ट करते
आणि संच उपपत्तीची स्वयंसिद्धे निवडण्यास आधार देते.
खालील युक्तिवाद हा रसेलचा संच का बनत नाही यासाठी आहे;
या मांडणीवर आधारित कोणतीही उपपत्ती \emph{सुसंगत} असते,
असे तो सिद्ध करत नाही. पूर्ण सुसंगततेसाठी
नेमकी स्वयंसिद्धे आणि स्वतंत्र युक्तिवाद आवश्यक असतो.

संचांची
संचयी-क्रमिक
संकल्पना
मानल्यास
संच
टप्प्याटप्प्याने
बनतात;
त्यांचे
!!{element}s
म्हणजे
\emph{आधीच}
उपलब्ध
असलेल्या
वस्तू असतात.
म्हणून
कोणत्याही
टप्प्या~$S$
साठी
पुढील संच
तयार करता
येतो:
\[
  R_S = \Setabs{x}{x \notin x \text{ आणि $x$ हा $S$ पूर्वी उपलब्ध होता}}
\]
रसेलची
विरोधापत्ती
सिद्ध करताना
वापरलेला
युक्तिवाद
दाखवतो
की~$R_S$
हा संच
स्वतः
टप्पा~$S$
पूर्वी
उपलब्ध
नसतो.
आणि
त्यात
व्याघात
नाही.
शिवाय
संचांची
संचयी-क्रमिक
संकल्पना
स्वीकारल्यास
रसेलचा
मूळ संच
तयार होईल
अशी
\emph{अपेक्षाच}
ठेवायला
नको होती.
कारण तो
‘‘कधीही
उपलब्ध
होतील’’
अशा
स्वतःचे
सदस्य
नसलेल्या
सर्व
संचांचा
संच
असता.
थोडक्यात,
संचयी-क्रमिक
मांडणी
मानल्यावर
रसेलचा
संच तयार
करता
येत
नाही
ही
(सिद्ध
होणारी)
गोष्ट
\emph{आश्चर्यकारक}
नाही;
याचाच
\emph{अंदाज}
आपण
केला
असता.

%In one sense, then, the cumulative-iterative conception of set yields a response to Russell's Paradox which is rather like the \emph{predicativist}'s response. After all, the predicativist said that the collection of all non-self-membered sets$_n$ is a set$_{n+1}$, and this set$_n$ will not be self-membered. (Indeed, most predicativists will treat it as \emph{ungrammatical} to try to ask whether a set is self-membered.)
%
%But there is an important difference between the cumulative-iterative approach and the predicativist approach: the cumulative-iterative approach treats all of our entities as being \emph{of the same kind}. The predicativist will presumably have an empty set$_0$ (i.e., a set$_0$ with no !!{element}s), and an empty set$_1$ (i.e., a set$_1$ with no !!{element}s), and these will be \emph{different entities}. But on the cumulative-iterative approach, there is just one empty set, $\emptyset$, and it will be available at every stage of the hierarchy.

%Indeed, the Axiom of Extensionality, stated at the start of this chapter, will hold true of the elements in this hierarchy (so that there can be \emph{only one} ``empty'' set). But I do not intend to use the cumulative-iterative conception to justify Extensionality (see \cite{Boolos1971}). Again: I suggest that we should accept Extensionality, just because we are interested in extensional collections.

\end{document}
